% GENERATED FILE -- do not edit by hand.
% Regenerate with: make index   (tools/gen_question_index.py)
\chapter{Question Index}
\label{chap:question_index}

\begin{tldr}
Every one of the 92 interview questions in Volume III, indexed by chapter, by round, and by frequency. This appendix is generated from the questions themselves, so it cannot drift from the chapters. Work the \freqhigh{} list first: those are the questions that come up in almost every loop.
\end{tldr}

\section{At a Glance}
\begin{table}[H]\centering
\caption{Question distribution}
\begin{tabular}{lr}
\toprule
\textbf{Category} & \textbf{Count} \\
\midrule
\quad \textbf{Total} & \textbf{92} \\
\quad L5 questions & 23 \\
\quad L6 questions & 53 \\
\quad L7 questions & 16 \\
\quad Breadth round & 26 \\
\quad Depth round & 42 \\
\quad Design round & 21 \\
\quad Behavioral round & 3 \\
\quad \freqhigh{} asked constantly & 29 \\
\quad \freqmed{} common & 51 \\
\quad \freqlow{} occasional & 12 \\
\bottomrule
\end{tabular}
\end{table}

\section{Start Here: The \freqhigh{} Questions}
These come up in almost every loop at this level. If you have one evening, these are the evening.
\begin{enumerate}[leftmargin=*]
\item \textbf{Ch.~1} \quad L5 \quad \textit{Conceptual} \quad Walk me through what happens, end to end, when a user types a query into a large-scale search engine
\item \textbf{Ch.~1} \quad L6 \quad \textit{System Design} \quad Design the query understanding system for an e-commerce search engine
\item \textbf{Ch.~1} \quad L6 \quad \textit{Debugging} \quad Your new L2 reranker improves offline NDCG by 4\%, but online conversions are flat. Walk through your diagnosis
\item \textbf{Ch.~2} \quad L5 \quad \textit{System Design} \quad Walk me through, mechanically, how BM25 over an inverted index returns the top-10 results from a billion documents in under 50\,ms
\item \textbf{Ch.~2} \quad L5 \quad \textit{Conceptual} \quad What do BM25's $k_1$ and $b$ actually control, and when would you change them from the defaults?
\item \textbf{Ch.~2} \quad L6 \quad \textit{Trade-off} \quad Lexical vs.\ dense vs.\ hybrid---argue the retrieval split for a marketplace search engine
\item \textbf{Ch.~3} \quad L5 \quad \textit{Conceptual} \quad How does HNSW work, and why is it fast?
\item \textbf{Ch.~3} \quad L6 \quad \textit{Architecture Design} \quad Index 500M $\times$ 768-d embeddings and serve top-100 under 15\,ms p99 on a machine with 64\,GB RAM---walk me through the design
\item \textbf{Ch.~4} \quad L5 \quad \textit{Conceptual} \quad Bi-encoder vs.\ cross-encoder: why not cross-encode everything, and what exactly does the cross-encoder buy?
\item \textbf{Ch.~4} \quad L6 \quad \textit{Debugging} \quad Your new reranker improved offline NDCG by 4\%, but online CTR dropped. Walk through the diagnosis
\item \textbf{Ch.~4} \quad L5 \quad \textit{Conceptual} \quad You run BM25 and dense retrieval in parallel. How do you combine them---and why does $0.5 \cdot \text{BM25} + 0.5 \cdot \text{cosine}$ fail?
\item \textbf{Ch.~4} \quad L6 \quad \textit{System Design} \quad Design semantic search over 100 million documents, end to end
\item \textbf{Ch.~5} \quad L6 \quad \textit{Communication} \quad Explain LambdaRank to a strong engineer who knows GBDT but has never done ranking
\item \textbf{Ch.~5} \quad L6 \quad \textit{Debugging} \quad Your click-trained ranker keeps favoring the items that have historically sat at position 1---better new items never rise. Diagnose and fix
\item \textbf{Ch.~5} \quad L6 \quad \textit{Trade-off} \quad GBDT or neural network for your L2 ranker? Decide for (a) a commerce search engine with rich engineered features, (b) a feed ranker over user history and item IDs
\item \textbf{Ch.~6} \quad L6 \quad \textit{Architecture Design} \quad Design a recommendation system for a video streaming platform with 100M users and 1M videos
\item \textbf{Ch.~6} \quad L5 \quad \textit{Conceptual} \quad How do you handle the cold start problem for new users and new items?
\item \textbf{Ch.~6} \quad L6 \quad \textit{Debugging} \quad Your recommendation model has great offline metrics but poor online performance. What went wrong?
\item \textbf{Ch.~6} \quad L6 \quad \textit{Architecture Design} \quad Design the ranking model for an e-commerce product ads system
\item \textbf{Ch.~6} \quad L6 \quad \textit{Estimation} \quad Estimate the memory required for DLRM embedding tables with 100M users, 10M items, and 1000 categorical features
\item \textbf{Ch.~6} \quad L6 \quad \textit{Debugging} \quad Your CTR model's offline AUC improved by 0.5\% but online revenue dropped 3\%. Diagnose
\item \textbf{Ch.~6} \quad L5 \quad \textit{Trade-off} \quad Two-Tower vs. interaction-based models for candidate generation---when would you use each?
\item \textbf{Ch.~6} \quad L5 \quad \textit{Conceptual} \quad Why do recommendation models need calibration? What happens if pCTR is systematically over-confident?
\item \textbf{Ch.~6} \quad L6 \quad \textit{Conceptual} \quad Where would you use an LLM in a recommendation stack today---and why won't it replace your ranker?
\item \textbf{Ch.~7} \quad L5 \quad \textit{Metric Derivation} \quad Derive NDCG from first principles, then compute NDCG@3 for a ranking with grades $(3, 0, 2)$, exponential gain, given the ideal available grades are $(3, 2, 0)$
\item \textbf{Ch.~7} \quad L6 \quad \textit{Debugging} \quad Your reranker improved offline NDCG@10 by 3\% on the judgment set, but the A/B test shows flat CTR. Walk me through your investigation
\item \textbf{Ch.~8} \quad L5 \quad \textit{System Design} \quad You own a search endpoint with a 250\,ms server-side p99 SLO. Walk me through the latency budget---where does the time go, and what do you cut when you're over?
\item \textbf{Ch.~8} \quad L6 \quad \textit{System Design} \quad Your search index must reflect catalog changes within 60 seconds---design it
\item \textbf{Ch.~8} \quad L6 \quad \textit{Debugging} \quad Relevance regressed after last week's embedder upgrade---find it
\end{enumerate}

\section{By Chapter}
\subsection*{Chapter 1: Search Systems and Query Understanding}
\begin{tabularx}{\textwidth}{@{}l l l X@{}}
\toprule
\textbf{Lvl} & \textbf{Freq} & \textbf{Type} & \textbf{Question} \\
\midrule
L5 & \freqhigh{} & Conceptual & Walk me through what happens, end to end, when a user types a query into a large-scale search engine \\
L6 & \freqhigh{} & System Design & Design the query understanding system for an e-commerce search engine \\
L6 & \freqmed{} & First Principles & Design spell correction for a commerce search engine. What breaks a naive dictionary approach? \\
L6 & \freqmed{} & Trade-off & How do you mine synonyms for query expansion without a hand-built ontology? \\
L7 & \freqmed{} & Trade-off & When should a query understanding signal change retrieval, and when should it only be a ranking feature? \\
L7 & \freqmed{} & System Design & You inherit a search system that is a single BM25 stage with a 15\% zero-result rate. Sequence your first three quarters of investment \\
L6 & \freqhigh{} & Debugging & Your new L2 reranker improves offline NDCG by 4\%, but online conversions are flat. Walk through your diagnosis \\
L5 & \freqmed{} & Conceptual & Explain the head/torso/tail decomposition of query traffic. Why must strategy differ by segment? \\
L6 & \freqmed{} & Trade-off & Where do LLMs fit in query understanding in 2026? Your PM wants ``LLM-powered search''---what do you actually build? \\
L5 & \freqlow{} & Debugging & Your search engine performs well in English but users report it is nearly useless in Japanese. Diagnose \\
L6 & \freqmed{} & Debugging & Search relevance metrics degraded starting last Tuesday. No model was retrained. Where do you look? \\
L6 & \freqmed{} & First Principles & Why is search a multi-stage funnel at all? Why not run your best model over the whole corpus for every query? \\
\bottomrule
\end{tabularx}

\subsection*{Chapter 2: Lexical Retrieval and Inverted Indexes}
\begin{tabularx}{\textwidth}{@{}l l l X@{}}
\toprule
\textbf{Lvl} & \textbf{Freq} & \textbf{Type} & \textbf{Question} \\
\midrule
L5 & \freqhigh{} & System Design & Walk me through, mechanically, how BM25 over an inverted index returns the top-10 results from a billion documents in under 50\,ms \\
L5 & \freqhigh{} & Conceptual & What do BM25's $k_1$ and $b$ actually control, and when would you change them from the defaults? \\
L6 & \freqmed{} & Debugging & A search for ``red dress'' ranks a product titled ``Red'' with ``dress'' scattered through a long description above an exact ``Red Dress'' title match. Diagnose and fix \\
L6 & \freqmed{} & Trade-off & Index-time vs.\ query-time synonyms---walk through the trade-offs and give your production policy \\
L6 & \freqmed{} & System Design & Explain WAND and block-max WAND. Why are they safe, and when does the pruning stop helping? \\
L5 & \freqlow{} & Conceptual & Why do search engines compress postings so aggressively when disk is cheap? Walk through the encoding stack \\
L7 & \freqmed{} & System Design & You must index 10B documents. Design the partitioning, tiering, and caching. Why did document partitioning beat term partitioning? \\
L6 & \freqmed{} & Conceptual & Explain SPLADE. What do the log-saturation and the FLOPS regularizer each do, and how does it compare to BM25 and dense retrieval at serving time? \\
L6 & \freqhigh{} & Trade-off & Lexical vs.\ dense vs.\ hybrid---argue the retrieval split for a marketplace search engine \\
L5 & \freqlow{} & Conceptual & How do phrase queries actually execute, and what do positional postings cost? When would you use phrase boosts vs.\ phrase matching vs.\ shingles? \\
\bottomrule
\end{tabularx}

\subsection*{Chapter 3: Vector Retrieval Theory}
\begin{tabularx}{\textwidth}{@{}l l l X@{}}
\toprule
\textbf{Lvl} & \textbf{Freq} & \textbf{Type} & \textbf{Question} \\
\midrule
L5 & \freqhigh{} & Conceptual & How does HNSW work, and why is it fast? \\
L6 & \freqhigh{} & Architecture Design & Index 500M $\times$ 768-d embeddings and serve top-100 under 15\,ms p99 on a machine with 64\,GB RAM---walk me through the design \\
L6 & \freqmed{} & Debugging & Your ANN recall dropped after a reindex---same code, same corpus size. Debug it \\
L5 & \freqmed{} & Conceptual & Why doesn't a $k$-d tree work for 768-dimensional embeddings, and what minimal changes rescue tree-based indexes? \\
L6 & \freqmed{} & Debugging & A recsys team L2-normalizes item embeddings from a dot-product-trained two-tower model so they can reuse a cosine HNSW index. Ranking quality drops. Diagnose and fix \\
L6 & \freqmed{} & Debugging & Your HNSW recall@10 drops from 0.95 to 0.6 when users apply a metadata filter matching 2\% of the corpus. Why, and what are the fixes? \\
L5 & \freqmed{} & Conceptual & Walk me through how a PQ index computes distances without decompressing anything, and where the time goes \\
L7 & \freqlow{} & Trade-off & When and why does ScaNN's anisotropic quantization beat plain PQ? \\
L6 & \freqlow{} & Estimation & You are offered an LSH family with $p_1 = 0.8$, $p_2 = 0.5$ at your target radius. Size an index for 100M vectors and decide whether to use it \\
L7 & \freqlow{} & Conceptual & Why does greedy search on a proximity graph find the nearest neighbor at all? What breaks in high dimensions, and which practical graph has worst-case guarantees? \\
L6 & \freqmed{} & Architecture Design & Tune an IVF index for 100M vectors: how many clusters? And what do you do when latency headroom remains but recall plateaus as you raise \texttt{nprobe}? \\
L6 & \freqlow{} & Debugging & Your ANN benchmark on synthetic Gaussian vectors shows terrible recall at any latency, but the same index on real embeddings is fine. Explain \\
\bottomrule
\end{tabularx}

\subsection*{Chapter 4: Neural Retrieval and Reranking}
\begin{tabularx}{\textwidth}{@{}l l l X@{}}
\toprule
\textbf{Lvl} & \textbf{Freq} & \textbf{Type} & \textbf{Question} \\
\midrule
L5 & \freqhigh{} & Conceptual & Bi-encoder vs.\ cross-encoder: why not cross-encode everything, and what exactly does the cross-encoder buy? \\
L6 & \freqmed{} & First Principles & Why do bi-encoders train with such large batches? Explain in-batch negatives and the logQ correction \\
L6 & \freqmed{} & Architecture Design & Design the training recipe for a domain bi-encoder from your search logs. Walk through negatives, denoising, and distillation \\
L6 & \freqmed{} & Trade-off & When does late interaction (ColBERT) earn its cost? Walk through the storage math \\
L7 & \freqmed{} & Trade-off & A listwise LLM reranker beats your production cross-encoder by 3 NDCG points offline---at 200$\times$ the cost. What ships? \\
L6 & \freqhigh{} & Debugging & Your new reranker improved offline NDCG by 4\%, but online CTR dropped. Walk through the diagnosis \\
L5 & \freqhigh{} & Conceptual & You run BM25 and dense retrieval in parallel. How do you combine them---and why does $0.5 \cdot \text{BM25} + 0.5 \cdot \text{cosine}$ fail? \\
L6 & \freqhigh{} & System Design & Design semantic search over 100 million documents, end to end \\
L6 & \freqmed{} & Debugging & Dense retrieval tanks on queries containing part numbers. Fix it without abandoning dense retrieval \\
L6 & \freqlow{} & Conceptual & Your reranker's score gates behavior downstream---``no results'' messaging and RAG context filtering. What breaks, and how do you make the scores trustworthy? \\
\bottomrule
\end{tabularx}

\subsection*{Chapter 5: Learning to Rank}
\begin{tabularx}{\textwidth}{@{}l l l X@{}}
\toprule
\textbf{Lvl} & \textbf{Freq} & \textbf{Type} & \textbf{Question} \\
\midrule
L5 & \freqmed{} & Conceptual & Compare pointwise, pairwise, and listwise learning to rank. When is each the right choice? \\
L6 & \freqhigh{} & Communication & Explain LambdaRank to a strong engineer who knows GBDT but has never done ranking \\
L6 & \freqhigh{} & Debugging & Your click-trained ranker keeps favoring the items that have historically sat at position 1---better new items never rise. Diagnose and fix \\
L5 & \freqmed{} & Applied & Walk me through constructing an LTR training set from a marketplace's search logs: labels, negatives, and splits \\
L5 & \freqmed{} & Trade-off & When would you insist on a pointwise objective even though the product is a ranked list? \\
L6 & \freqhigh{} & Trade-off & GBDT or neural network for your L2 ranker? Decide for (a) a commerce search engine with rich engineered features, (b) a feed ranker over user history and item IDs \\
L6 & \freqmed{} & Applied & How would you estimate position-bias propensities without degrading the user experience? \\
L6 & \freqmed{} & Debugging & Your new LTR model improves overall NDCG but tail-query relevance regresses. Why, and what do you change? \\
L6 & \freqmed{} & Judgment Call & Editorial judgments say document A beats B for this query; click data says B massively outperforms A. Which do you trust, and what do you do? \\
L7 & \freqmed{} & System Design & Design the end-to-end unbiased-LTR loop for a search product---logging through training through evaluation---and tell me where it silently breaks \\
L7 & \freqlow{} & Depth & Your IPS-weighted training runs are unstable---a few examples dominate the gradient and validation NDCG oscillates. What is happening and what are your options? \\
\bottomrule
\end{tabularx}

\subsection*{Chapter 6: Recommendation Systems}
\begin{tabularx}{\textwidth}{@{}l l l X@{}}
\toprule
\textbf{Lvl} & \textbf{Freq} & \textbf{Type} & \textbf{Question} \\
\midrule
L6 & \freqhigh{} & Architecture Design & Design a recommendation system for a video streaming platform with 100M users and 1M videos \\
L5 & \freqhigh{} & Conceptual & How do you handle the cold start problem for new users and new items? \\
L6 & \freqhigh{} & Debugging & Your recommendation model has great offline metrics but poor online performance. What went wrong? \\
L6 & \freqhigh{} & Architecture Design & Design the ranking model for an e-commerce product ads system \\
L7 & \freqmed{} & Architecture Design & Design a real-time personalized notification system \\
L7 & \freqlow{} & Architecture Design & Design the embedding infrastructure for a recommendation system with 1B items \\
L6 & \freqhigh{} & Estimation & Estimate the memory required for DLRM embedding tables with 100M users, 10M items, and 1000 categorical features \\
L6 & \freqmed{} & Estimation & Your ads ranking model needs to score 100 candidates in $<$50ms. What architecture constraints does this impose? \\
L7 & \freqlow{} & Estimation & Calculate the daily training data volume for an ads system serving 1B impressions/day \\
L6 & \freqhigh{} & Debugging & Your CTR model's offline AUC improved by 0.5\% but online revenue dropped 3\%. Diagnose \\
L7 & \freqmed{} & Debugging & Your recommendation model shows high engagement but users are churning. What is happening? \\
L5 & \freqhigh{} & Trade-off & Two-Tower vs. interaction-based models for candidate generation---when would you use each? \\
L6 & \freqmed{} & Trade-off & Shared-bottom vs. MMoE vs. PLE for multi-task ads modeling---give me a decision framework \\
L6 & \freqmed{} & Conceptual & Explain the sample selection bias problem in CVR prediction. How does ESMM solve it? \\
L5 & \freqhigh{} & Conceptual & Why do recommendation models need calibration? What happens if pCTR is systematically over-confident? \\
L6 & \freqmed{} & Trade-off & You run retrieval for a marketplace with 200M listings and heavy daily churn. Two-tower + ANN or generative retrieval with semantic IDs? \\
L6 & \freqhigh{} & Conceptual & Where would you use an LLM in a recommendation stack today---and why won't it replace your ranker? \\
\bottomrule
\end{tabularx}

\subsection*{Chapter 7: Search and RecSys Evaluation}
\begin{tabularx}{\textwidth}{@{}l l l X@{}}
\toprule
\textbf{Lvl} & \textbf{Freq} & \textbf{Type} & \textbf{Question} \\
\midrule
L5 & \freqhigh{} & Metric Derivation & Derive NDCG from first principles, then compute NDCG@3 for a ranking with grades $(3, 0, 2)$, exponential gain, given the ideal available grades are $(3, 2, 0)$ \\
L5 & \freqmed{} & Conceptual & You have binary judgments only. When do Precision@$k$, Recall@$k$, MRR, and MAP each answer the right question---and construct a case where two of them disagree about which of two systems is better \\
L6 & \freqhigh{} & Debugging & Your reranker improved offline NDCG@10 by 3\% on the judgment set, but the A/B test shows flat CTR. Walk me through your investigation \\
L6 & \freqmed{} & Trade-off & Interleaving vs.\ A/B testing for a ranking change---how does team-draft interleaving work, why is it more sensitive, and when would it mislead you? \\
L6 & \freqmed{} & Debugging & You replace a lexical retriever with a dense retriever. Offline NDCG on the existing judgment set drops. Is the new system worse? \\
L6 & \freqmed{} & Protocol Design & A colleague evaluates a new sequential recommender with a random 80/20 interaction split and HR@10 against 100 sampled negatives, and reports a 12\% win. What is wrong, and what protocol do you require before believing it? \\
L6 & \freqmed{} & Judgment Call & Your new recommender lifts NDCG and short-term engagement, but intra-list diversity drops, catalog coverage falls, and exposure Gini rises 6 points. Do you ship it---and how should the evaluation have been set up so this is not a debate? \\
L7 & \freqmed{} & First Principles & Without launching it, estimate what CTR a new ranking policy would achieve, using only logs from the current system. Derive the estimator, its requirements, and its failure modes \\
L6 & \freqmed{} & Program Design & You want to replace most crowd relevance labeling with an LLM judge. Design the program so the labels are trustworthy---and tell me what stays human \\
L7 & \freqlow{} & Program Design & You own relevance for a search org of several teams. Design the evaluation program: what gets measured, at what layer, on what cadence---and how does a change get to ship? \\
\bottomrule
\end{tabularx}

\subsection*{Chapter 8: Production Retrieval and RAG Integration}
\begin{tabularx}{\textwidth}{@{}l l l X@{}}
\toprule
\textbf{Lvl} & \textbf{Freq} & \textbf{Type} & \textbf{Question} \\
\midrule
L5 & \freqhigh{} & System Design & You own a search endpoint with a 250\,ms server-side p99 SLO. Walk me through the latency budget---where does the time go, and what do you cut when you're over? \\
L6 & \freqhigh{} & System Design & Your search index must reflect catalog changes within 60 seconds---design it \\
L5 & \freqmed{} & Conceptual & Why are deletes and updates hard in HNSW, and what does your vector database actually do when you call \texttt{delete()}? \\
L6 & \freqhigh{} & Debugging & Relevance regressed after last week's embedder upgrade---find it \\
L6 & \freqmed{} & System Design & You have no labels in production. How do you know your retrieval quality hasn't regressed---and how do you find out fast? \\
L6 & \freqmed{} & Trade-off & Traffic to your LLM-answer product is expensive. Design a semantic cache---and tell me how it goes wrong \\
L6 & \freqmed{} & Trade-off & You're building the index behind a RAG product over 100M documents. Chunk-level or doc-level vectors---and what does the choice do to your index? \\
L5 & \freqmed{} & Trade-off & We have 20M vectors, moderate QPS, and we already run Postgres and OpenSearch. Do we need a vector database? \\
L7 & \freqmed{} & System Design & Cut retrieval serving cost 5$\times$ without visible quality loss \\
L7 & \freqmed{} & Estimation & A new embedding model tops the leaderboard. Do you re-embed your 100M-document corpus? Walk through the decision and the cost \\
\bottomrule
\end{tabularx}

\section{By Round}
The rounds a loop is actually made of. Drill one column at a time.
\subsection*{Breadth (26 questions)}
\begin{itemize}[leftmargin=*]
\item \textbf{Ch.~1} \quad L5 \quad Walk me through what happens, end to end, when a user types a query into a large-scale search engine
\item \textbf{Ch.~1} \quad L6 \quad Design spell correction for a commerce search engine. What breaks a naive dictionary approach?
\item \textbf{Ch.~1} \quad L5 \quad Explain the head/torso/tail decomposition of query traffic. Why must strategy differ by segment?
\item \textbf{Ch.~1} \quad L6 \quad Why is search a multi-stage funnel at all? Why not run your best model over the whole corpus for every query?
\item \textbf{Ch.~2} \quad L5 \quad What do BM25's $k_1$ and $b$ actually control, and when would you change them from the defaults?
\item \textbf{Ch.~2} \quad L5 \quad Why do search engines compress postings so aggressively when disk is cheap? Walk through the encoding stack
\item \textbf{Ch.~2} \quad L6 \quad Explain SPLADE. What do the log-saturation and the FLOPS regularizer each do, and how does it compare to BM25 and dense retrieval at serving time?
\item \textbf{Ch.~2} \quad L5 \quad How do phrase queries actually execute, and what do positional postings cost? When would you use phrase boosts vs.\ phrase matching vs.\ shingles?
\item \textbf{Ch.~3} \quad L5 \quad How does HNSW work, and why is it fast?
\item \textbf{Ch.~3} \quad L5 \quad Why doesn't a $k$-d tree work for 768-dimensional embeddings, and what minimal changes rescue tree-based indexes?
\item \textbf{Ch.~3} \quad L5 \quad Walk me through how a PQ index computes distances without decompressing anything, and where the time goes
\item \textbf{Ch.~3} \quad L7 \quad Why does greedy search on a proximity graph find the nearest neighbor at all? What breaks in high dimensions, and which practical graph has worst-case guarantees?
\item \textbf{Ch.~4} \quad L5 \quad Bi-encoder vs.\ cross-encoder: why not cross-encode everything, and what exactly does the cross-encoder buy?
\item \textbf{Ch.~4} \quad L6 \quad Why do bi-encoders train with such large batches? Explain in-batch negatives and the logQ correction
\item \textbf{Ch.~4} \quad L5 \quad You run BM25 and dense retrieval in parallel. How do you combine them---and why does $0.5 \cdot \text{BM25} + 0.5 \cdot \text{cosine}$ fail?
\item \textbf{Ch.~4} \quad L6 \quad Your reranker's score gates behavior downstream---``no results'' messaging and RAG context filtering. What breaks, and how do you make the scores trustworthy?
\item \textbf{Ch.~5} \quad L5 \quad Compare pointwise, pairwise, and listwise learning to rank. When is each the right choice?
\item \textbf{Ch.~5} \quad L5 \quad Walk me through constructing an LTR training set from a marketplace's search logs: labels, negatives, and splits
\item \textbf{Ch.~5} \quad L6 \quad How would you estimate position-bias propensities without degrading the user experience?
\item \textbf{Ch.~6} \quad L5 \quad How do you handle the cold start problem for new users and new items?
\item \textbf{Ch.~6} \quad L6 \quad Explain the sample selection bias problem in CVR prediction. How does ESMM solve it?
\item \textbf{Ch.~6} \quad L5 \quad Why do recommendation models need calibration? What happens if pCTR is systematically over-confident?
\item \textbf{Ch.~6} \quad L6 \quad Where would you use an LLM in a recommendation stack today---and why won't it replace your ranker?
\item \textbf{Ch.~7} \quad L5 \quad You have binary judgments only. When do Precision@$k$, Recall@$k$, MRR, and MAP each answer the right question---and construct a case where two of them disagree about which of two systems is better
\item \textbf{Ch.~7} \quad L7 \quad Without launching it, estimate what CTR a new ranking policy would achieve, using only logs from the current system. Derive the estimator, its requirements, and its failure modes
\item \textbf{Ch.~8} \quad L5 \quad Why are deletes and updates hard in HNSW, and what does your vector database actually do when you call \texttt{delete()}?
\end{itemize}

\subsection*{Depth (42 questions)}
\begin{itemize}[leftmargin=*]
\item \textbf{Ch.~1} \quad L6 \quad How do you mine synonyms for query expansion without a hand-built ontology?
\item \textbf{Ch.~1} \quad L7 \quad When should a query understanding signal change retrieval, and when should it only be a ranking feature?
\item \textbf{Ch.~1} \quad L6 \quad Your new L2 reranker improves offline NDCG by 4\%, but online conversions are flat. Walk through your diagnosis
\item \textbf{Ch.~1} \quad L6 \quad Where do LLMs fit in query understanding in 2026? Your PM wants ``LLM-powered search''---what do you actually build?
\item \textbf{Ch.~1} \quad L5 \quad Your search engine performs well in English but users report it is nearly useless in Japanese. Diagnose
\item \textbf{Ch.~1} \quad L6 \quad Search relevance metrics degraded starting last Tuesday. No model was retrained. Where do you look?
\item \textbf{Ch.~2} \quad L6 \quad A search for ``red dress'' ranks a product titled ``Red'' with ``dress'' scattered through a long description above an exact ``Red Dress'' title match. Diagnose and fix
\item \textbf{Ch.~2} \quad L6 \quad Index-time vs.\ query-time synonyms---walk through the trade-offs and give your production policy
\item \textbf{Ch.~2} \quad L6 \quad Lexical vs.\ dense vs.\ hybrid---argue the retrieval split for a marketplace search engine
\item \textbf{Ch.~3} \quad L6 \quad Your ANN recall dropped after a reindex---same code, same corpus size. Debug it
\item \textbf{Ch.~3} \quad L6 \quad A recsys team L2-normalizes item embeddings from a dot-product-trained two-tower model so they can reuse a cosine HNSW index. Ranking quality drops. Diagnose and fix
\item \textbf{Ch.~3} \quad L6 \quad Your HNSW recall@10 drops from 0.95 to 0.6 when users apply a metadata filter matching 2\% of the corpus. Why, and what are the fixes?
\item \textbf{Ch.~3} \quad L7 \quad When and why does ScaNN's anisotropic quantization beat plain PQ?
\item \textbf{Ch.~3} \quad L6 \quad You are offered an LSH family with $p_1 = 0.8$, $p_2 = 0.5$ at your target radius. Size an index for 100M vectors and decide whether to use it
\item \textbf{Ch.~3} \quad L6 \quad Your ANN benchmark on synthetic Gaussian vectors shows terrible recall at any latency, but the same index on real embeddings is fine. Explain
\item \textbf{Ch.~4} \quad L6 \quad When does late interaction (ColBERT) earn its cost? Walk through the storage math
\item \textbf{Ch.~4} \quad L7 \quad A listwise LLM reranker beats your production cross-encoder by 3 NDCG points offline---at 200$\times$ the cost. What ships?
\item \textbf{Ch.~4} \quad L6 \quad Your new reranker improved offline NDCG by 4\%, but online CTR dropped. Walk through the diagnosis
\item \textbf{Ch.~4} \quad L6 \quad Dense retrieval tanks on queries containing part numbers. Fix it without abandoning dense retrieval
\item \textbf{Ch.~5} \quad L6 \quad Your click-trained ranker keeps favoring the items that have historically sat at position 1---better new items never rise. Diagnose and fix
\item \textbf{Ch.~5} \quad L5 \quad When would you insist on a pointwise objective even though the product is a ranked list?
\item \textbf{Ch.~5} \quad L6 \quad GBDT or neural network for your L2 ranker? Decide for (a) a commerce search engine with rich engineered features, (b) a feed ranker over user history and item IDs
\item \textbf{Ch.~5} \quad L6 \quad Your new LTR model improves overall NDCG but tail-query relevance regresses. Why, and what do you change?
\item \textbf{Ch.~5} \quad L7 \quad Your IPS-weighted training runs are unstable---a few examples dominate the gradient and validation NDCG oscillates. What is happening and what are your options?
\item \textbf{Ch.~6} \quad L6 \quad Your recommendation model has great offline metrics but poor online performance. What went wrong?
\item \textbf{Ch.~6} \quad L6 \quad Estimate the memory required for DLRM embedding tables with 100M users, 10M items, and 1000 categorical features
\item \textbf{Ch.~6} \quad L6 \quad Your ads ranking model needs to score 100 candidates in $<$50ms. What architecture constraints does this impose?
\item \textbf{Ch.~6} \quad L7 \quad Calculate the daily training data volume for an ads system serving 1B impressions/day
\item \textbf{Ch.~6} \quad L6 \quad Your CTR model's offline AUC improved by 0.5\% but online revenue dropped 3\%. Diagnose
\item \textbf{Ch.~6} \quad L7 \quad Your recommendation model shows high engagement but users are churning. What is happening?
\item \textbf{Ch.~6} \quad L5 \quad Two-Tower vs. interaction-based models for candidate generation---when would you use each?
\item \textbf{Ch.~6} \quad L6 \quad Shared-bottom vs. MMoE vs. PLE for multi-task ads modeling---give me a decision framework
\item \textbf{Ch.~6} \quad L6 \quad You run retrieval for a marketplace with 200M listings and heavy daily churn. Two-tower + ANN or generative retrieval with semantic IDs?
\item \textbf{Ch.~7} \quad L5 \quad Derive NDCG from first principles, then compute NDCG@3 for a ranking with grades $(3, 0, 2)$, exponential gain, given the ideal available grades are $(3, 2, 0)$
\item \textbf{Ch.~7} \quad L6 \quad Your reranker improved offline NDCG@10 by 3\% on the judgment set, but the A/B test shows flat CTR. Walk me through your investigation
\item \textbf{Ch.~7} \quad L6 \quad Interleaving vs.\ A/B testing for a ranking change---how does team-draft interleaving work, why is it more sensitive, and when would it mislead you?
\item \textbf{Ch.~7} \quad L6 \quad You replace a lexical retriever with a dense retriever. Offline NDCG on the existing judgment set drops. Is the new system worse?
\item \textbf{Ch.~8} \quad L6 \quad Relevance regressed after last week's embedder upgrade---find it
\item \textbf{Ch.~8} \quad L6 \quad Traffic to your LLM-answer product is expensive. Design a semantic cache---and tell me how it goes wrong
\item \textbf{Ch.~8} \quad L6 \quad You're building the index behind a RAG product over 100M documents. Chunk-level or doc-level vectors---and what does the choice do to your index?
\item \textbf{Ch.~8} \quad L5 \quad We have 20M vectors, moderate QPS, and we already run Postgres and OpenSearch. Do we need a vector database?
\item \textbf{Ch.~8} \quad L7 \quad A new embedding model tops the leaderboard. Do you re-embed your 100M-document corpus? Walk through the decision and the cost
\end{itemize}

\subsection*{Design (21 questions)}
\begin{itemize}[leftmargin=*]
\item \textbf{Ch.~1} \quad L6 \quad Design the query understanding system for an e-commerce search engine
\item \textbf{Ch.~1} \quad L7 \quad You inherit a search system that is a single BM25 stage with a 15\% zero-result rate. Sequence your first three quarters of investment
\item \textbf{Ch.~2} \quad L5 \quad Walk me through, mechanically, how BM25 over an inverted index returns the top-10 results from a billion documents in under 50\,ms
\item \textbf{Ch.~2} \quad L6 \quad Explain WAND and block-max WAND. Why are they safe, and when does the pruning stop helping?
\item \textbf{Ch.~2} \quad L7 \quad You must index 10B documents. Design the partitioning, tiering, and caching. Why did document partitioning beat term partitioning?
\item \textbf{Ch.~3} \quad L6 \quad Index 500M $\times$ 768-d embeddings and serve top-100 under 15\,ms p99 on a machine with 64\,GB RAM---walk me through the design
\item \textbf{Ch.~3} \quad L6 \quad Tune an IVF index for 100M vectors: how many clusters? And what do you do when latency headroom remains but recall plateaus as you raise \texttt{nprobe}?
\item \textbf{Ch.~4} \quad L6 \quad Design the training recipe for a domain bi-encoder from your search logs. Walk through negatives, denoising, and distillation
\item \textbf{Ch.~4} \quad L6 \quad Design semantic search over 100 million documents, end to end
\item \textbf{Ch.~5} \quad L7 \quad Design the end-to-end unbiased-LTR loop for a search product---logging through training through evaluation---and tell me where it silently breaks
\item \textbf{Ch.~6} \quad L6 \quad Design a recommendation system for a video streaming platform with 100M users and 1M videos
\item \textbf{Ch.~6} \quad L6 \quad Design the ranking model for an e-commerce product ads system
\item \textbf{Ch.~6} \quad L7 \quad Design a real-time personalized notification system
\item \textbf{Ch.~6} \quad L7 \quad Design the embedding infrastructure for a recommendation system with 1B items
\item \textbf{Ch.~7} \quad L6 \quad A colleague evaluates a new sequential recommender with a random 80/20 interaction split and HR@10 against 100 sampled negatives, and reports a 12\% win. What is wrong, and what protocol do you require before believing it?
\item \textbf{Ch.~7} \quad L6 \quad You want to replace most crowd relevance labeling with an LLM judge. Design the program so the labels are trustworthy---and tell me what stays human
\item \textbf{Ch.~7} \quad L7 \quad You own relevance for a search org of several teams. Design the evaluation program: what gets measured, at what layer, on what cadence---and how does a change get to ship?
\item \textbf{Ch.~8} \quad L5 \quad You own a search endpoint with a 250\,ms server-side p99 SLO. Walk me through the latency budget---where does the time go, and what do you cut when you're over?
\item \textbf{Ch.~8} \quad L6 \quad Your search index must reflect catalog changes within 60 seconds---design it
\item \textbf{Ch.~8} \quad L6 \quad You have no labels in production. How do you know your retrieval quality hasn't regressed---and how do you find out fast?
\item \textbf{Ch.~8} \quad L7 \quad Cut retrieval serving cost 5$\times$ without visible quality loss
\end{itemize}

\subsection*{Behavioral (3 questions)}
\begin{itemize}[leftmargin=*]
\item \textbf{Ch.~5} \quad L6 \quad Explain LambdaRank to a strong engineer who knows GBDT but has never done ranking
\item \textbf{Ch.~5} \quad L6 \quad Editorial judgments say document A beats B for this query; click data says B massively outperforms A. Which do you trust, and what do you do?
\item \textbf{Ch.~7} \quad L6 \quad Your new recommender lifts NDCG and short-term engagement, but intra-list diversity drops, catalog coverage falls, and exposure Gini rises 6 points. Do you ship it---and how should the evaluation have been set up so this is not a debate?
\end{itemize}

